% te amo mi vida <3  

\documentclass[8pt,a4paper,landscape]{extarticle} 

% Lade dein eigenes Cheatsheet-Paket!
\usepackage{TeXothek_Template}




%-------------------------------------------
%%%%%%  CHEATSHEET STARTS HERE  %%%%%%%%%%%%
%-------------------------------------------
\begin{document}
\tiny\color{MyGray}
\setlength{\lineskiplimit}{-\maxdimen}
\linespread{1.2}\selectfont
\begin{multicols*}{6}




%---------First Row-------------------------
{\huge\color{MyPurple1}\textbf{Stochastische Signale}}
\vspace{-5pt}
\section*{Wahrscheinlichkeitsräume:}

\text{Wahrscheinlichkeitsraum $(\Omega, \mathbb{F}, \mathbb{P})$ besteht aus:}
  
  \CheatsheetHeading{Ergebnismenge $\Omega$: }{}

  $\Omega = \{\omega_1, \omega_2, \dots\}$: \text{Menge aller möglichen Ergebnisse.}
%
  \CheatsheetHeading{Ereignisalgebra $\mathbb{F}$:}{}

  $\mathbb{F} = \{A_1, A_2, \dots\}$: \text{Menge von Ereignissen} $A_i \subseteq \Omega$.
%
    \CheatsheetItemListStart
      \CheatsheetItem{Eigenschaften: }{}{}
    \CheatsheetItemListEnd
%
      $\sbullet \ \Omega \in \mathbb{F} $ \\

      $\sbullet A_i \in \mathbb{F} \ \Rightarrow \ A_i^c \in \mathbb{F} $ \\
      $\sbullet A_1, \dots, A_k \in \mathbb{F} \ \Rightarrow \ \cup_{i=1}^{k} A_i \in \mathbb{F} $ \\
%
      \text{- Impliziert: }$\sbullet\emptyset \in \mathbb{F} \text{ }(\Omega^c = \emptyset)\text{ }$ 
      $\sbullet A_i\backslash A_k \in \mathbb{F}\text{ }$
      $\sbullet \cap_{i = 1}^k \in \mathbb{F}$
%
    \CheatsheetItemListStart
     \CheatsheetItem{Größe: }{}{}
    \CheatsheetItemListEnd 
%
      $\sbullet|\mathbb{F}| \hat{=} 2^{\text{Anzahl an disjunkter Teilmengen}}$ 
%
  \CheatsheetHeading{Wahrscheinlichkeitsmaß P:}{}

  \text{Ordnet jedem Ereignis A eine Wahrscheinlichkeit zu.}
%
    \CheatsheetItemListStart
        \CheatsheetItem{Axiome von Kolmogorow:}{}{}
    \CheatsheetItemListEnd
%    
      $\sbullet \mathbb{P}(A) \ge 0 \ \Rightarrow \ \mathbb{P}: \mathbb{F} \mapsto [0,1] \\$ 
      $\sbullet \mathbb{P}(\Omega) = 1$ \\
      $\sbullet \mathbb{P}(\cup_{i = 1}^{\infty} A_i) = \sum_{i = 1}^{\infty} \mathbb{P}(A_i) \ \text{wenn: } \overbrace{A_i \cap A_j = \emptyset, \forall i \neq j}^{\text{(disjunkt)}}$
%
  \CheatsheetHeading{Grundlegende Rechenregeln:}{}

    $\sbullet \mathbb{P}(A^c) = 1 - \mathbb{P}(A) $ \\
    $\sbullet \mathbb{P}(\emptyset) = 0 $ \hspace{3pt} $\sbullet \mathbb{P}(\Omega) = 1 $ \\
    $\sbullet \mathbb{P}(A \setminus B) = \mathbb{P}(A) - \mathbb{P}(A \cap B) $ \\
    $\sbullet \mathbb{P}(A \cup B) = \mathbb{P}(A) + \mathbb{P}(B) - \mathbb{P}(A \cap B) $ \\
    $\sbullet \mathbb{P}(A \cap B) = \mathbb{P}(A) + \mathbb{P}(B) - \mathbb{P}(A \cup B) $ \\
    $\sbullet A \subset B \Rightarrow \mathbb{P}(A) \leq \mathbb{P}(B) $

\section{Bedingte Wahrscheinlichkeit:}

$\hat{=}$ Der Wahrscheinlichkeit für ein Ereignis A, unter der Bedingung, 
dass bereits ein anderes Ereignis B eingetreten ist. \\
\par\vspace*{-3pt}
$\color{orange}\sbullet \mathbb{P}_B(A) = \mathbb{P}(A | B) = \frac{\mathbb{P}(A \cap B)}{\mathbb{P}(B)}$
\vspace{3pt}
%
  \CheatsheetHeading{Multiplikationssatz und dessen Erweiterung:}{}

    \vspace{-1pt}
    \fboxit{\mathbb{P}(A \cap B) = \mathbb{P}(A | B) \cdot \mathbb{P}(B) = \mathbb{P}(B | A) \cdot \mathbb{P}(A) }{}
    \vspace{-2pt}
    \fboxit{\mathbb{P}(A \cap B \cap C) = \mathbb{P}(A|B \cap C) \cdot \mathbb{P}(B \cap C) = \\
    \text{}\hspace{37pt}\,= \mathbb{P}(A|B \cap C) \cdot \mathbb{P}(B|C) \cdot \mathbb{P}(C)}{}
    \vspace{0pt}

\section{Totale Wahrscheinlichkeit:}

\vspace{-5pt}
\fimagetext{0.5}{Images/Totale_Wahrscheinlichkeit}{$\Omega = \cup_{i = 1}^{n} B_i \ $ \\
$ \text{mit:} \ B_i \cap B_j = \emptyset $ \\
\text{} \\
$\mathbb{P}(A \cap B) = $ \\
$ = \mathbb{P}(A | B_i) \cdot \mathbb{P}(B_i)$
} 
%
\par\vspace*{10pt}
{$\color{orange}\sbullet \mathbb{P}(A) = \sum_{i=1}^{n} \mathbb{P}(A | B_i) \cdot \mathbb{P}(B_i)$}
%
  \CheatsheetHeading{Satz von Bayes:}{}
    \vspace{-2pt}
    \fboxit{\mathbb{P}(B_j|A) = \frac{\mathbb{P}(A|B_j) \cdot \mathbb{P}(B_j)}{\sum_{i=1}^{n} \mathbb{P}(A|B_i) \cdot \mathbb{P}(B_i)}}{}
    \vspace{5pt}
    \fboxit{\mathbb{P}(B_j|A) = \frac{\mathbb{P}(A|B_j) \cdot \mathbb{P}(B_j)}{\mathbb{P}(A)}}{}
    \vspace{0pt}

\section{Stochastische Unabhängigkeit:}
%
Zwei Ereignisse A und B sind stochastisch unabhängig, wenn das Eintreten von
A die Wahrscheinlichkeit von B nicht beeinflusst.
%
\par
$\color{orange}\sbullet \mathbb{P}(A \cap B) = \mathbb{P}(A) \cdot \mathbb{P}(B) \Rightarrow \mathbb{P}(A | B) = \mathbb{P}(A) $

\section{Zufallsvariablen}

%
 \CheatsheetItemListStart
  \CheatsheetItem{Definition: }{$A \subset \Omega \ \text{oder } A \in \mathbb{F} \\
    X:\Omega \rightarrow \Omega ^{\backprime} \text{, } \omega \rightarrow \text{X(w) heißt Zufallsvariable, wenn} \\
    \text{für jedes }  A^{\backprime} \in \mathbb{F}^{\backprime} \text{ ein } A \in \mathbb{F}
    \text{ existiert, sodass } \\
    A = \{\omega \in  \Omega | X(\omega) \in A^{\backprime}\}$}{}
  \vspace*{-13pt}
  \CheatsheetItem{\color{MyGreen} Beispiel: }{2-mal Würfeln und Augenzahl addieren}{}
 \CheatsheetItemListEnd
%
\par\vspace{10pt}
\fimagetext{1}{images/Zufallsvariablen.png}{}
\par\vspace*{7pt}

        
















%---------Second Row------------------------
\columnbreak

\section{Wahrscheinlichkeitsverteilungen}
%
  \CheatsheetHeading{Kumulative Verteilungsfunktion: }{(CDF bzw. KVF)}

  $ F_x(x) = \mathbb{P}({X \leq x}) \text{; } X \text{: Zufallsvariable} \text{; } x \text{: Realisierung}$ 
  \vspace*{8pt}
  \fimagetext{1}{images/CDFs.png}{}
  \vspace*{-2pt}
%
  \CheatsheetItemListStart
    \CheatsheetItem{Eigenschaften:}{}{}
  \CheatsheetItemListEnd
%
    $\sbullet F_x(x) \text{ monoton wachsend und rechtsseitig stetig} $ \\
    $\sbullet \lim\limits_{x \to \infty} F_x(x) = 0 \text{ ; } \lim\limits_{x \to \infty} F_x(x) = 1 $ \\
    \par\vspace*{-5pt}
    $\sbullet F_x(x) \ge 0 $
%
  \par\vspace*{-6pt}
  \fsubcols{\CheatsheetSubHeading{Diskrete ZV:}{}
    $\sum_{\mathbb{R}} p_X(x) = 1$
%
    \CheatsheetItemListStart
      \color{MyPurple2}\CheatsheetItem{Wahrs.massenfkt:}{ PMF}{}
    \CheatsheetItemListEnd
%
    $\sbullet p_X(x) = \mathbb{P}(\{X = x\})$
%
    \CheatsheetItemListStart
      \color{MyPurple2}\CheatsheetItem{Kumula. Vert.fkt:}{ CDF}{}
    \CheatsheetItemListEnd
%
    $\sbullet F_X(x) = \sum_{k \le x} p_X(k)$

  }
  {\CheatsheetSubHeading{Stetige ZV:}{}
%
    \vspace*{0pt}
    \CheatsheetItemListStart
      \color{MyPurple2}\CheatsheetItem{Wahrs.dichtefkt:}{ PDF}{}
    \CheatsheetItemListEnd
%
    $\sbullet f_X(x) = \frac{d}{dx} F_X(x)$
%
    \CheatsheetItemListStart
      \color{MyPurple2}\CheatsheetItem{Kumula. Vert.fkt:}{ CDF}{}
    \CheatsheetItemListEnd
%
    $\sbullet F_X(x) = \int_{-\infty}^{x} f_X(\xi) \, d\xi$
%
  
  }
%
    \vspace*{40pt}
    \CheatsheetItemListStart
      \color{orange}\CheatsheetItem{Anderes zu Stetigen ZV:}{}{}
    \CheatsheetItemListEnd
%
    $\sbullet F_X(b) - F_X(a) = \int_{a}^{b} f_X(x) \, dx$ \\
    $\sbullet \mathbb{P}(\{X = a\}) = \int_{a}^{a} f_X(x) \, dx = 0 \quad \sbullet \int_{-\infty}^{\infty} f_X(x) \, dx = 1$
%

\section{Bedingte Zufallsvariablen}
%
\CheatsheetItemListStart
  \CheatsheetItem{Gegebenes Ereignis A: }{$\quad F_{X|A}(x|A) = \mathbb{P}(\{X \le x|A\})$}{}
\CheatsheetItemListEnd
%
\CheatsheetItemListStart
  \CheatsheetItem{Gegebene ZV Y: }{}{}
\CheatsheetItemListEnd
%
\vspace*{5pt}
\faligntext{
  \sbullet F_{X|Y}(x|y) &= \mathbb{P}(\{X \le x|Y = y\}) \\[7pt]
  \sbullet p_{X|Y}(x|y) &= \frac{p_{X,Y}(x,y)}{p_Y(y)} \quad \sbullet
  f_{X|Y}(x|y) = \frac{f_{X,Y}(x,y)}{f_Y(y)} \\
}{}
%

\vspace*{4pt}

\section{Mehrdim. Zufallsvariablen}
%
Gegeben sind n Zufallsvariablen $X_1$,$X_2$, ... ,$X_n$
%
\CheatsheetHeading{Mehrdimensionale CDF (KVF)}{}
%
\par
$\sbullet F_{X_1,..,X_n}(x_1,.., x_n) = \mathbb{P}(\{X_1 \le x_1, X_2 \le x_2,.., X_n \le x_n\})$ \\
$\sbullet F_{X_1,X_2}(x_1, x_2) = \mathbb{P}(A) = \int_{-\infty}^{x_1} \int_{-\infty}^{x_2} f_{X_1,X_2}(\xi, \eta) \, d\eta \, d\xi$ \\
%
\vspace*{17pt}
\fimagetext{0.5}{Images/CDF_moredim.png}{}
\vspace*{7pt}
%
  \CheatsheetHeading{diskret (PMF):}{}
% 
  \par
  $p_{X_1, \dots, X_n}(x_1, \dots, x_n) = \mathbb{P}(\{X_1 = x_1, \dots, X_n = x_n\})$
%
  \CheatsheetHeading{stetig (PDF):}{}
%
\par
$f_{X_1,\dots,X_n}(x_1, \dots, x_n) = \frac{\partial^n}{\partial x_1, \dots, \partial x_n} F_{X_1,\dots,X_n}(x_1, \dots, x_n)$ \\[2pt]
$F_{X_1,.,X_n}(x_1,.., x_n) = \int_{-\infty}^{x_n} .. \int_{-\infty}^{x_1} f_{X_1,..,X_n}(\xi_1,.., \xi_n) \, d\xi_1 .. d\xi_n $ \\[2pt]
$F_{X,Y} = F_{Y,X}$

\CheatsheetItemListStart
  \CheatsheetItem{\color{MyGreen} Beispiel:}{}{}
\CheatsheetItemListEnd
%
$f_{x,y}(x,y) = \frac{4}{19}(1 + xy)$ \quad ($x \in [2,3] \wedge  y \in [1,2]$) \\
$\text{ }\hspace{22pt} = 0$ \quad (sonst)
\vspace*{20pt}
\fimagetext{0.5}{Images/CDF_moredim_example.png}{}
%
\par\vspace*{35pt}
  \begingroup
    \setlength{\fboxsep}{1.5pt}
      \begin{varwidth}{\dimexpr\linewidth-5\fboxsep\relax} %sets maximum width on coloum width
        \ensuremath{   
        F_{X,Y}(x,y) = \begin{cases}    
          \quad 0 \quad \quad ; x,y \in \mathrm{I, IV, VII, VIII, IX} \\
          \frac{4}{19} \int_2^x \int_1^y (1 + \xi\eta) \, d\eta \, d\xi & ; x,y \in \mathrm{V} \\
          \quad 1 & ; x,y \in \mathrm{III} \\
          \frac{4}{19} \int_2^3 \int_1^y (1 + \xi\eta) \, d\eta \, d\xi & ; x,y \in \mathrm{VI} \\
          \frac{4}{19} \int_2^x \int_1^2 (1 + \xi\eta) \, d\eta \, d\xi & ; x,y \in \mathrm{II}
          \end{cases}
        }
      \end{varwidth}
    \vspace{-3pt}\raggedright\tiny
  \endgroup
  \par\vspace*{2pt}\noindent
  \vspace{12pt}

\section{Marginalisierung}

%
Prinzip: Lasse alle vernachlässigbaren ZV gegen unendlich gehen.
%
\CheatsheetHeading{Kumulative Verteilungsfkt. (KVF):}{}
%
\par
%
$F_X(x) = \mathbb{P}(\{X \le x\}) = \\
\mathbb{P}(\{X \le x, Y\}) = \lim_{y \to \infty} F_{X,Y}(x, y) = F_{X,Y}(x, \infty)$ \\
$p_{X_1}(x_1) = \sum_{x_2, \dots, x_n} p_{X_1, \dots, X_n}(x_1, \dots, x_n)$ \\
$f_{X_1}(x_1) = \int_{-\infty}^{\infty} \dots \int_{-\infty}^{\infty} f_{X_1, \dots, X_n}(x_1, \dots, x_n) \, dx_2 \dots dx_n$ \\
 


%---------Third Row-------------------------
\columnbreak
%

\section{Unabhängigkeit von ZV}
%
Zufallsvariablen $X_1, ... ,X_i$ sind stochastisch unabhängig, wenn \\ 
$\sbullet  \mathbb{P}(\{X_1 \le x_1, \dots, X_n \le x_n\}) = \prod_{i=1}^{n} \mathbb{P}(\{X_i \le x_i\})$ \\
und somit \\
$\sbullet F_{X_1, \dots, X_n}(x_1, \dots, x_n) = \prod_{i=1}^{n} F_{X_i}(x_i)$ \\
$\sbullet p_{X_1, \dots, X_n}(x_1, \dots, x_n) = \prod_{i=1}^{n} p_{X_i}(x_i)$ \\
$\sbullet f_{X_1, \dots, X_n}(x_1, \dots, x_n) = \prod_{i=1}^{n} f_{X_i}(x_i)$

\section{Funktionen von Zufallsvariablen}
{\color{gray} $Y = g(X) \rightarrow $ X u. Y sind nicht sto. unabhängig} 
%
\CheatsheetHeading{Transformation mit monoton und differenziebaren Fkt:}{}
%
\par
\vspace{18pt}
\fimagetext{0.37}{Images/ZV_Trafo.png}{$\vspace{2pt} \break f_Y(y) = \left. \frac{f_X(x)}{|g'(x)|} \right|_{x=g^{-1}(y)} \vspace{3pt}\break \text{(} \Rightarrow \text{mit Umkehrfunktion: } g^{-1}(y)\text{)}$}
\vspace{6pt}
%
\CheatsheetHeading{Transformation mit nur differenziebaren Funktionen}{}
\par
{\color{gray} Strategie: Aufteilen in monotone Funktionen.} \\
%
$\text{Allgemein: } f_Y(y) = \sum_{i=1}^{N} \frac{f_X(x_i)}{\left| \frac{dg(x)}{dx} \right|_{x=x_i}} \qquad i \in \{1, \dots, N\}$ \\
%
\par
\vspace{-2pt}
$(x_i \text{ sind Nst. von } y - g(x) = 0)$
\CheatsheetItemListStart
 \CheatsheetItem{\color{MyGreen} Beispiel:}{}{}
\CheatsheetItemListEnd
\par
\vspace{40pt}
\fimagetext{0.4}{Images/ZV_Trafo2.png}{Gegeben: \break $\sbullet$ Wertebereich: $\quad  y \in [0; 1]$ $\sbullet \ g'(x) = 2x$ \break \text{Daraus folgt:} \break $\Rightarrow \quad y = x^2 \rightarrow x_{1/2} = \pm\sqrt{y}$ 
$\Rightarrow \quad f_Y(y) = \vspace{2pt} \break = \frac{f_X(x_1)}{|g'(x_1)|} + \frac{f_X(x_2)}{|g'(x_2)|} = \vspace{2pt} \break = \frac{\frac{1}{2}}{|2\sqrt{y}|} + \frac{\frac{1}{2}}{|-2\sqrt{y}|} = \frac{1}{2\sqrt{y}} $}
\vspace{25pt}
%
\CheatsheetItemListStart
 \CheatsheetItem{\color{MyGreen} Beispiel für das Interval einer ZV-Trafo:}{}{}
\CheatsheetItemListEnd
%
\par\vspace*{5pt}
  \begingroup
    \setlength{\fboxsep}{1.5pt}
      \begin{varwidth}{\dimexpr\linewidth-5\fboxsep\relax} %sets maximum width on coloum width
        \ensuremath{   
          f_X(x) = \begin{cases} 
          \frac{1}{\sqrt{2}} & 0 < x < \sqrt{2} \\
          0 & \text{sonst} 
          \end{cases} 
        }
      \end{varwidth}
    \vspace{-3pt}\raggedright\tiny
  \endgroup
  \par\vspace*{2pt}\noindent
  \vspace{2pt}
% 
\par\vspace*{5pt}
  \begingroup
    \setlength{\fboxsep}{1.5pt}
      \begin{varwidth}{\dimexpr\linewidth-5\fboxsep\relax} %sets maximum width on coloum width
        \ensuremath{   
g(x) = \begin{cases} 
x^2 & x \leq \sqrt{2} \\ 
2 & x > \sqrt{2} 
\end{cases}
       }
      \end{varwidth}
    \vspace{-3pt}\raggedright\tiny
  \endgroup
  \par\vspace*{2pt}\noindent
  \vspace{2pt}
%
\break
$Y = g(X) \rightarrow [0, 2]$
%
\CheatsheetHeading{Lineare Trafo. von ZV:}{}
%
\par
$\text{Sei: } Y = aX + b \iff g(x) = ax + b \implies f_Y(y) = \frac{1}{|a|} f_X\left(\frac{y-b}{a}\right)$

\par\vspace*{8pt}
  \begingroup
    \setlength{\fboxsep}{1.5pt}
      \begin{varwidth}{\dimexpr\linewidth-5\fboxsep\relax} %sets maximum width on coloum width
        \ensuremath{   
        F_Y(y) = \begin{cases} 
        F_X\left(\frac{y-b}{a}\right) & a > 0 \\ 
        1 - F_X\left(\frac{y-b}{a}\right) & a < 0 
        \end{cases}
        }
      \end{varwidth}
    \vspace{-3pt}\raggedright\tiny
  \endgroup
  \par\vspace*{2pt}\noindent
  \vspace{2pt}
%
\CheatsheetHeading{Summe unabhängiger ZV:}{}
%
\vspace{8pt}
\fimagetext{0.6}{Images/ZV_Trafo3.png}{}
\vspace{8pt}
\par
$\text{X \& Y sind stochastisch unabhängig:} \quad  Z = X + Y \break
F_Z(z) = \mathbb{P}(\{Z \leq z\}) \break
\textcolor{orange}{\implies f_{Z=X+Y}(z) = (f_X * f_Y)(z) = \int_{-\infty}^{\infty} f_X(z - y) f_Y \, dy}$
%

\section{Stochastische standardmodelle}
%
\vspace{-5pt}
\fsubcols{\CheatsheetSubHeading{Diskrete Modelle}{}
\par
$\sbullet$ Diskrete Gleichverteilung \\
$\sbullet$ Geometrische Verteilung \\
$\sbullet$ Bernoulliverteilung \\
$\sbullet$ Binominalverteilung \\
$\sbullet$ Poisson-Verteilung \\
}
{\CheatsheetSubHeading{stetige Modelle}{}
\par
$\sbullet$ Gleichverteilung \\
$\sbullet$ Exponentialverteilung \\
$\sbullet$ Gauß/Normalverteilung \\
}
%
\vspace{35pt}
\CheatsheetHeading{Gedächtnislosigkeit:}{}
\par
ZV ist gedächtnislos, falls \\ {\color{orange} $\mathbb{P}(\{X > a + b\} | \{x > a\}) = \mathbb{P}(\{X > b\})$}








%---------Fourth Row------------------------
\columnbreak
\vspace*{-10pt}
\par\noindent 
\CheatsheetHeading{Gleichverteilung:}{stetig/diskret}
%
\vspace*{6pt}
\par
  \begingroup
    \setlength{\fboxsep}{1.5pt}
      \begin{varwidth}{\dimexpr\linewidth-5\fboxsep\relax} 
        \ensuremath{   
        \begin{array}{lll}
        \hspace*{-5pt}
        \text{\textbf{\textit{diskret:}}} \ 
        p_X(x) = \frac{1}{|\Omega|} \ 
        \text{(Münze, Würfel)} \\
        
        \hspace*{-5pt}
        \text{\textbf{\textit{stetig:}}} \ 
        f_X(x) = \begin{cases} 
        \frac{1}{b-a} \text{, } x \in [a, b] \\ 
        0 \text{, sonst} 
        \end{cases}   
        
        \hspace*{-3pt}
        F_X(x) = \begin{cases} 
        0 & \text{, } x < a \\ 
        \frac{x-a}{b-a} & \text{, } x \in [a, b] \\ 
        1 & \text{, } x > b 
        \end{cases}
        \end{array}
        }
      \end{varwidth}
    \vspace{-3pt}\raggedright\tiny
  \endgroup
  \par\noindent
  \vspace{25pt}
%
\par
\fimagetext{1}{Images/Gleichverteilung.png}{}
\vspace{-3pt}
%
%
\fsubcols
{\CheatsheetItemListStart
\CheatsheetItem{Erwartungswert:}{}{}
\CheatsheetItemListEnd
$\mathbb{E}[X] = \frac{a+b}{2}$
}
{
\CheatsheetItemListStart
\CheatsheetItem{Varianz:}{}{}
\CheatsheetItemListEnd
$\operatorname{Var}[X] = \frac{(b-a)^2}{12}$
}
\vspace{4pt}
\CheatsheetItemListStart
\CheatsheetItem{Charackeristische Fkt:}{}{}
\CheatsheetItemListEnd
$\varphi_X(\omega) = \frac{\mathrm{e}^{\mathrm{j}\omega b} - \mathrm{e}^{\mathrm{j}\omega a}}{\mathrm{j}\omega(b-a)}$
%
\vspace{2pt}
\CheatsheetHeading{Bernoulliverteilung:}{diskret}

\vspace{10pt}
\begin{varwidth}{\dimexpr\linewidth-5\fboxsep\relax}
  \ensuremath{
    \begin{array}{lll}
      \hspace*{-5pt}
      p_X(x) = \begin{cases}
        1-p & ,\ x = 0 \\
        p   & ,\ x = 1 \\
        0   & ,\ \text{sonst}
      \end{cases} 
      \hspace*{7pt}
      F_X(x) = \begin{cases}
        0      & ,\ x < 0 \\
        1-p    & ,\ 0 \le x < 1 \\
        1      & ,\ x \ge 1
      \end{cases}
    \end{array}
  }
\end{varwidth}
\vspace{10pt}

\vspace{14pt}
\fimagetext{0.3}{Images/Bernoulliverteilung.png}{\text{2 Ereignisse: Erfolg/Miserflog}}
\vspace{2pt}
\vspace{-3pt}
\fsubcols
{\CheatsheetItemListStart
\CheatsheetItem{Erwartungswert:}{}{}
\CheatsheetItemListEnd
$\mathbb{E}[X] = p$
}
{
\CheatsheetItemListStart
\CheatsheetItem{Varianz:}{}{}
\CheatsheetItemListEnd
$\operatorname{Var}[X] = p(1-p)$
}
\vspace{4pt}
\CheatsheetItemListStart
\CheatsheetItem{Charakteristische Fkt:}{}{}
\CheatsheetItemListEnd
$\varphi_X(\omega) = \mathbb{E}[\mathrm{e}^{\mathrm{j}\omega X}] = (1-p) + p \,\mathrm{e}^{\mathrm{j}\omega}$


\CheatsheetHeading{Binominalverteilung:}{diskret}

\vspace{6pt}
\begin{varwidth}{\dimexpr\linewidth-5\fboxsep\relax}
  \ensuremath{
    \begin{array}{lll}
      \hspace*{-5pt}
      p_X(k) = \begin{cases}
        {n \choose k} p^k (1-p)^{n-k} & ,\ k = 0,1,\ldots,n \\
        0 & ,\ \text{sonst}
      \end{cases} \\ 
      \end{array}
  }
\end{varwidth}
\vspace{8pt}
\par
$F_X(k) = \sum_{i=0}^{\lfloor k\rfloor} {n \choose i} p^i (1-p)^{n-i}$
    

\vspace{26pt}
\fimagetext{1}{Images/Binomialverteilung.png}{}
\vspace{2pt}
\fsubcols
{\CheatsheetItemListStart
\CheatsheetItem{Erwartungswert:}{}{}
\CheatsheetItemListEnd
$\mathbb{E}[X] = np$
}
{
\CheatsheetItemListStart
\CheatsheetItem{Varianz:}{}{}
\CheatsheetItemListEnd
$\operatorname{Var}[X] = np(1-p)$
}
\vspace{4pt}
\CheatsheetItemListStart
\CheatsheetItem{Charakteristische Fkt:}{}{}
\CheatsheetItemListEnd
$\varphi_X(\omega) = \bigl((1-p) + p \,\mathrm{e}^{\mathrm{j}\omega}\bigr)^n$

\CheatsheetHeading{Poisson-Verteilung:}{diskret}

\vspace{10pt}
\begin{varwidth}{\dimexpr\linewidth-5\fboxsep\relax}
  \ensuremath{
    \begin{array}{lll}
      \hspace*{-5pt}
      p_X(k) = \begin{cases}
        \mathrm{e}^{-\lambda} \frac{\lambda^k}{k!} & ,\ k = 0,1,2,\ldots \\
        0 & ,\ \text{sonst}
      \end{cases} \\
      \hspace*{7pt}
    \end{array}
  }
\end{varwidth}
\vspace{5pt}

\text{k: Anzahl der Erfolge, p: Erfolgswahrscheinlichkeit}
\text{Folge von n Bernoulli-Experimenten.} \\
$F_X(k) = \sum_{i=0}^{\lfloor k\rfloor} \mathrm{e}^{-\lambda} \frac{\lambda^i}{i!}$
\vspace{26pt}
\fimagetext{1}{Images/Binomialverteilung.png}{}
\vspace{5pt}
\vspace{-3pt}
\fsubcols
{\CheatsheetItemListStart
\CheatsheetItem{Erwartungswert:}{}{}
\CheatsheetItemListEnd
$\mathbb{E}[X] = \lambda$
}
{
\CheatsheetItemListStart
\CheatsheetItem{Varianz:}{}{}
\CheatsheetItemListEnd
$\operatorname{Var}[X] = \lambda$
}
\vspace{4pt}
\CheatsheetItemListStart
\CheatsheetItem{Charakteristische Fkt:}{}{}
\CheatsheetItemListEnd
$\varphi_X(\omega) = \exp\left(\lambda \left(\mathrm{e}^{\mathrm{j}\omega} - 1\right)\right)$




\columnbreak

\vspace*{-10pt}


\CheatsheetHeading{Geometrische Verteilung:}{diskret/gedächtnislos}

\vspace{20pt}
\begin{varwidth}{\dimexpr\linewidth-5\fboxsep\relax}
  \ensuremath{
    \begin{array}{lll}
      \hspace*{-5pt}
      p_X(k) = \begin{cases}
        p\,(1-p)^{\,k-1} & ,\ k = 1,2,3,\dots \\
        0 & ,\ \text{sonst}
      \end{cases} \\[8pt]
      \hspace*{-5pt}
      F_X(k) = \begin{cases}
        0 & ,\ k < 1 \\
        1 - (1-p)^k & ,\ k = 1,2,3,\dots \\
        1 & ,\ \text{sonst}
      \end{cases}
    \end{array}
  }
\end{varwidth}
\vspace{10pt}

\vspace{38pt}
\fimagetext{1}{Images/Geometrische_Verteilung.png}{}

\vspace{12pt}
\fsubcols
{\CheatsheetItemListStart
\CheatsheetItem{Erwartungswert:}{}{}
\CheatsheetItemListEnd
$\mathbb{E}[X] = \frac{1}{p}$
}
{
\CheatsheetItemListStart
\CheatsheetItem{Varianz:}{}{}
\CheatsheetItemListEnd
$\operatorname{Var}[X] = \frac{1-p}{p^{2}}$
}
\vspace{4pt}
\CheatsheetItemListStart
\CheatsheetItem{Charakteristische Fkt:}{}{}
\CheatsheetItemListEnd
$\varphi_X(\omega) = \frac{p\,\mathrm{e}^{\mathrm{j}\omega}}{\,1 - (1-p)\,\mathrm{e}^{\mathrm{j}\omega}\,}$

\vspace{2pt}
\CheatsheetHeading{Exponentialverteilung:}{stetig/gedächtnislos}

\vspace{20pt}
\begin{varwidth}{\dimexpr\linewidth-5\fboxsep\relax}
  \ensuremath{
    \begin{array}{lll}
      \hspace*{-5pt}
      f_X(x) = \begin{cases}
        \lambda\,\mathrm{e}^{-\lambda x} & ,\ x \ge 0 \\
        0 & ,\ \text{sonst}
      \end{cases} \\[8pt]
      \hspace*{-5pt}
      F_X(x) = \begin{cases}
        0 & ,\ x < 0 \\
        1 - \mathrm{e}^{-\lambda x} & ,\ x \ge 0
      \end{cases}
    \end{array}
  }
\end{varwidth}
\vspace{10pt}

\vspace{34pt}
\fimagetext{1}{Images/Exponentialverteilung.png}{}

\vspace{12pt}
\fsubcols
{\CheatsheetItemListStart
\CheatsheetItem{Erwartungswert:}{}{}
\CheatsheetItemListEnd
$\mathbb{E}[X] = \frac{1}{\lambda}$
}
{
\CheatsheetItemListStart
\CheatsheetItem{Varianz:}{}{}
\CheatsheetItemListEnd
$\operatorname{Var}[X] = \frac{1}{\lambda^{2}}$
}
\vspace{4pt}
\CheatsheetItemListStart
\CheatsheetItem{Charakteristische Fkt:}{}{}
\CheatsheetItemListEnd
$\varphi_X(\omega) = \frac{\lambda}{\lambda - \mathrm{j}\omega}$

\vspace{2pt}
\CheatsheetHeading{Normalverteilung:}{stetig}

\vspace{10pt}
\begin{varwidth}{\dimexpr\linewidth-5\fboxsep\relax}
  \ensuremath{
    \begin{array}{lll}
      \hspace*{-5pt}
      f_X(x) = \frac{1}{\sqrt{2\pi\sigma^2}} \, \mathrm{e}^{-\frac{(x-\mu)^2}{2\sigma^2}} \\[6pt]
      \hspace*{-5pt}
      F_X(x) = \frac{1}{2} \left[ 1 + \operatorname{erf}\left( \frac{x-\mu}{\sigma\sqrt{2}} \right) \right] = \Phi\left(\frac{x-\mu}{\sigma}\right)
    \end{array}
  }
\end{varwidth}
\vspace{10pt}

\vspace{22pt}
\fimagetext{1}{Images/Normalverteilung.png}{}

\vspace{10pt}
\fsubcols
{\CheatsheetItemListStart
\CheatsheetItem{Erwartungswert:}{}{}
\CheatsheetItemListEnd
$\mathbb{E}[X] = \mu$
}
{
\CheatsheetItemListStart
\CheatsheetItem{Varianz:}{}{}
\CheatsheetItemListEnd
$\operatorname{Var}[X] = \sigma^2$
}
\vspace{4pt}
\CheatsheetItemListStart
\CheatsheetItem{Charakteristische Fkt:}{}{}
\CheatsheetItemListEnd
$\varphi_X(\omega) = \mathrm{e}^{\mathrm{j}\omega\mu - \frac{1}{2}\sigma^2\omega^2}$

\CheatsheetItemListStart
\CheatsheetItem{Eigenschaften: }{symmetrisch, bei $x=\mu$ maximal}{}
\CheatsheetItemListEnd
$X \sim \mathcal{N}(\mu, \sigma)$










%---------Fifth Row-------------------------
\section{Erwartungswert \& Varianz}

\CheatsheetHeading{Erwartungswert:}{ mittlerer Wert einer ZV, \quad  $\mathbb{E}[X] = \mu_X$}

      $\mathbb{E}[X] = \sum_{x \in \Omega'} x \cdot p_X(x) \quad \text{diskrete Zufallsvariable}$ \\[8pt]
      $\mathbb{E}[X] = \underbrace{\int\limits_{-\infty}^{\infty} x \cdot f_X(x) \, dx}_{\text{1. Moment}} \quad \text{stetige Zufallsvariable}$
\vspace{15pt}

\CheatsheetItemListStart
  \CheatsheetItem{Eigenschaften:}{}{}
\CheatsheetItemListEnd
\par\vspace{2pt}
Linearität: $\mathbb{E}[\alpha X + \beta Y] = \alpha \mathbb{E}[X] + \beta \mathbb{E}[Y] = \alpha\mu + \beta\eta$ \\
Monotonie: $X \le Y \Rightarrow \mathbb{E}[X] \le \mathbb{E}[Y]$ \\
$\sbullet \mathbb{E}[a] = a$ \\
$\sbullet \mathbb{E}[-X] = -\mathbb{E}[X]$ \\
$\sbullet \mathbb{E}[X + \beta] = \mathbb{E}[X] + \beta$ \\
$\sbullet$ \text{bei einer Fkt: } $Y = g(x)$ vereinfachen durch einsetzen! \\
\hspace*{4pt} $ \mathbb{E}[X Y] = \iint\limits_{\mathbb{R}^2} x y f_{XY}(x,y) \, dx \, dy$ \\

$\sbullet$ \text{sto. unabhängig, dann:} \quad
$\mathbb{E}[X Y] = \mathbb{E}[X]\mathbb{E}[Y]$ \\
{\color{red}\textbf{!} $\mathbb{E}[X]\mathbb{E}[Y] = \mathbb{E}[XY] \nRightarrow \text{unabh.}$ \textbf{!}}
\CheatsheetHeading{Varianz:}{Maß fur die Abweichung vom Erwartungswert}

$\operatorname{Var}[X] = \mathbb{E}[(X-\mathbb{E}[X])^2] = \mathbb{E}[X^2] - (\mathbb{E}[X])^2$

\CheatsheetItemListStart
  \CheatsheetItem{Eigenschaften:}{}{}
\CheatsheetItemListEnd
$\sbullet \operatorname{Var}[-X] = \operatorname{Var}[X]$ \\
$\sbullet \operatorname{Var}[aX] = a^2\operatorname{Var}[X],\ a \in \mathbb{R}$ \\
$\sbullet \operatorname{Var}[a] = 0$ \\
$\sbullet \operatorname{Var}[aX+\beta] = a^2\operatorname{Var}[X]$ \\
$\sbullet \operatorname{Var}[X+Y] = \operatorname{Var}[X] + \operatorname{Var}[Y] + 2\operatorname{Cov}[X,Y]$ \\
$\sbullet \operatorname{Var}[X-Y] = \operatorname{Var}[X] + \operatorname{Var}[Y] - 2\operatorname{Cov}[X,Y]$

\CheatsheetHeading{Kovarianz:}{linearer Zusammenhang zwischen zwei ZV}

$\operatorname{Cov}[X,Y] = \mathbb{E}[(X-\mathbb{E}[X])(Y-\mathbb{E}[Y])]$ \\
$\operatorname{Cov}[X,Y] = \mathbb{E}[XY]-\mathbb{E}[X]\mathbb{E}[Y]$

\CheatsheetItemListStart
  \CheatsheetItem{Eigenschaften:}{}{}
\CheatsheetItemListEnd
$\sbullet \operatorname{Cov}[X,X] = \operatorname{Var}[X]$ \\
$\sbullet \operatorname{Cov}[X,Y] = \operatorname{Cov}[Y,X]$ \\
$\sbullet \operatorname{Cov}[aX+\beta,Y] = a\operatorname{Cov}[X,Y]$ \\
$\sbullet X \perp Y \Rightarrow \operatorname{Cov}[X,Y] = 0$ \quad (Umkehrung i.A. falsch)

\CheatsheetHeading{Korrelationskoeffizient:}{normierte Kovarianz}

$\rho_{X,Y} = \frac{\operatorname{Cov}[X,Y]}{\sqrt{\operatorname{Var}[X] \operatorname{Var}[Y]}}$ 
\vspace*{2pt}

Es gilt: $\rho_{X,Y} \in [-1,1]$ \\
$\sbullet \rho_{X,Y} < 0$ negativ korreliert, \\
$\sbullet \rho_{X,Y}=0$ unkorreliert, \\
$\sbullet \rho_{X,Y} > 0$ positiv korreliert.

\CheatsheetItemListStart
  \CheatsheetItem{Unkorreliert vs. orthogonal:}{}{}
\CheatsheetItemListEnd
Unkorreliert: $\operatorname{Cov}[X,Y] = 0$ \\
Orthogonal: $\mathbb{E}[XY] = 0$


\section{Erzeugend/Charakteristisch Fkt.}

\CheatsheetHeading{Wahrscheinlichkeits-erzeugende Funktion:}{}
\vspace*{-4pt}
\fboxit{G_X(z) = \mathbb{E}[z^X] = \sum_{k=0}^{\infty} z^k p_X(k), \quad z \in \mathbb{C},\ |z| \le 1}{} \\
\vspace*{-7pt}
\CheatsheetItemListStart
  \CheatsheetItem{Eigenschaften:}{}{}
\CheatsheetItemListEnd
\vspace*{0pt}
$\sbullet \mathbb{P}(X=n) = \frac{1}{n!}\left.\frac{d^n}{dz^n}G_X(z)\right|_{z=0}$ \\ [3pt]
$\sbullet \mathbb{E}[X] = \left.\frac{d}{dz}G_X(z)\right|_{z=1}$ \\ [3pt]
$\sbullet \mathbb{E}[X^2]-\mathbb{E}[X] = \left.\frac{d^2}{dz^2}G_X(z)\right|_{z=1}$ \\ [3pt]
$\sbullet \operatorname{Var}[X] = \left.\frac{d^2}{dz^2}G_X(z)\right|_{z=1} - \mathbb{E}[X]^2 + \mathbb{E}[X]$ \\ [3pt]
$\sbullet$ Wenn nur $G_X(z)$ fur eine diskrete Zufallsvariable bekannt ist, konnen daraus Verteilung, $\mathbb{E}[X]$ und $\operatorname{Var}[X]$ bestimmt werden. \\
\vspace*{-7pt}

\CheatsheetHeading{Momentenerzeugende Funktion:}{}

\vspace*{3pt}
\fboxit{M_X(s) = \mathbb{E}[\mathrm{e}^{sX}] = \int_{-\infty}^{\infty} f_X(x)\mathrm{e}^{sx}\,dx = \sum_{k=0}^{\infty} \frac{s^k}{k!} \mathbb{E}[X^k]}{} \\

\vspace*{-7pt}
\CheatsheetItemListStart
  \CheatsheetItem{Eigenschaft:}{}{}
\CheatsheetItemListEnd
$\sbullet \mathbb{E}[X^n] = \left.\frac{d^n}{ds^n}M_X(s)\right|_{s=0}$

\CheatsheetHeading{Charakteristische Funktion:}{stetige ZV}
\vspace*{-6pt}
\fboxit{\varphi_X(\omega) = \mathbb{E}[\mathrm{e}^{\mathrm{j}\omega X}] = \int_{-\infty}^{\infty} f_X(x)\mathrm{e}^{\mathrm{j}\omega x}\,dx}{} \\
\vspace*{-9pt}
\CheatsheetItemListStart
  \CheatsheetItem{Eigenschaft:}{}{}
\CheatsheetItemListEnd
$\sbullet \mathbb{E}[X^n] = \left.\frac{1}{\mathrm{j}^n}\frac{d^n}{d\omega^n}\varphi_X(\omega)\right|_{\omega=0}$

\CheatsheetHeading{Summe unabhangiger ZV:}{}

Sind $X,Y$ unabhangig und $Z=X+Y$, dann gilt:

$\sbullet G_Z(z)=G_X(z)\cdot G_Y(z)$ \\
$\sbullet M_Z(s)=M_X(s)\cdot M_Y(s)$ \\
$\sbullet \varphi_Z(\omega)=\varphi_X(\omega)\cdot \varphi_Y(\omega)$

\CheatsheetHeading{Zentraler Grenzwertsatz:}{}

Seien $X_i$ i.i.d. mit $\mathbb{E}[X_i]=\mu<\infty$ und $\operatorname{Var}[X_i]=\sigma^2<\infty$, dann:
\vspace*{-5pt} \par
\fboxit{Z_n = \frac{\sum_{i=1}^{n}X_i - n\mu}{\sigma\sqrt{n}} \xrightarrow{d} \mathcal{N}(0,1)}{} \\
Dann konvergiert die Verteilung der standardisierten Summe gegen die Standardnormalverteilung. \\
$\Rightarrow \mathbb{E}[Z_n]=0,\ \operatorname{Var}[Z_n]=1\ \text{(im Grenzfall)}$ \\
d.h. Die Summe ist (asymptotisch) gaußverteilt, unabhangig von der Verteilung der Summanden.


\columnbreak

\section{Zufallsfolgen}

\CheatsheetHeading{Definition:}{}

$X_i:\Omega \to \mathbb{R},\ i\in\mathbb{N}$

Bei $N\to\infty$ wäre die vollstandige Beschreibung über die gemeinsame CDF
$F_{X_1,\dots,X_N}(x_1,\dots,x_N)$ sehr komplex.

\CheatsheetHeading{Vereinfachungsschritt:}{Momente 1./2. Ordnung}

Betrachte Erwartungswertfolge, Varianzfolge und \\
(Auto-)Korrelation/Kovarianz:

$\sbullet\ \mu_X(i)=\mathbb{E}[X_i]$ ({\color{orange}Erwartungswertfolge}) \\
$\sbullet\ \sigma_X^2(i)=\operatorname{Var}[X_i]=\mathbb{E}[X_i^2]-\mathbb{E}[X_i]^2$ ({\color{orange}Varianzfolge}) \\
$\sbullet\ r_X(k,l)=\mathbb{E}[X_kX_l]$ ({\color{orange}Autokorrelationsfolge}) \\
$\sbullet\ c_X(k,l)=\operatorname{Cov}[X_k,X_l]=r_X(k,l)-\mu_X(k)\mu_X(l)$ \\
$\phantom{\sbullet\ }=\mathbb{E}[(X_k-\mathbb{E}[X_k])(X_l-\mathbb{E}[X_l])]$ ({\color{orange}Autokovarianzfolge})

\CheatsheetHeading{Zwei Zufallsfolgen $X_i,Y_i$:}{}

$\sbullet\ r_{XY}(k,l)=\mathbb{E}[X_kY_l]$ ({\color{orange}Kreuzkorrelationsfolge}) \\
$\sbullet\ c_{XY}(k,l)=\operatorname{Cov}[X_k,Y_l]=r_{XY}(k,l)-\mu_X(k)\mu_Y(l)$ \\
$\phantom{\sbullet }=\mathbb{E}[(X_k-\mathbb{E}[X_k])(Y_l-\mathbb{E}[Y_l])]$ ({\color{orange}Kreuzkovarianzfolge}) \\
\textbf{Orthogonal:} $r_{XY}(k,l)=0$ \quad \textbf{Unkorreliert:} $c_{XY}(k,l)=0$

\CheatsheetHeading{i.i.d.-Folge:}{unabhangig und identisch verteilt}

$X_i$ i.i.d. $\Rightarrow \mathbb{E}[X_i]=\mu,\ \operatorname{Var}[X_i]=\sigma^2,\ \operatorname{Cov}[X_i,X_j]=0\ (i\neq j)$

\CheatsheetHeading{Stationaritat:}{}

{\color{orange}Streng stationar:} \\
$F_{X_{i_1},\dots,X_{i_m}}(x_1,\dots,x_m)=F_{X_{i_1+k},\dots,X_{i_m+k}}(x_1,\dots,x_m)$ \\
{\color{orange}WSS:} \\
$\mu_X(i)=\mu$ konstant,\ $r_X(k,l)=r_X(k-l)$,\\ 
$c_X(k,l)=c_X(k-l)$ 
\CheatsheetItemListStart
  \CheatsheetItem{wichtige Relationen:}{}{}
\CheatsheetItemListEnd
stationär $\Rightarrow$ WSS \\
WSS $\nRightarrow$ Stationär 

\CheatsheetHeading{Random Walk:}{}

{\color{orange}Definition:} mathematisches Modell für eine Folge zufälliger Schritte,   
um zufallige Prozesse/Bewegungen zu modellieren.

{\color{MyBlue}Grundkonzept:} Start bei einem Punkt; pro Schritt entweder vorwärts $(+\delta)$
mit Wkt. $p$ oder ruckwarts $(-\delta)$ mit Wkt. $1-p$.

{\color{orange}Mathematisch:} $n\in\mathbb{N}$ Schritte, $X_i\in\{+\delta,-\delta\}$ (i.i.d.) \\
$\sbullet S_n=\sum_{i=1}^{n}X_i$, für kleine und große n nicht stationär \\
$\sbullet \mathbb{P}(\{X_i=+\delta\})=p,\quad \sbullet \mathbb{P}(\{X_i=-\delta\})=1-p$ \\
$\sbullet \mathbb{E}[X_i]=(2p-1)\delta,\quad \sbullet \operatorname{Var}[X_i]=4p(1-p)\delta^2$ \\
$\sbullet \mu_S(n)=\mathbb{E}[S_n]=n(2p-1)\delta$ \\
$\sbullet \sigma_S^2(n)=\operatorname{Var}[S_n]=4np(1-p)\delta^2$ \\
Fur $p=\frac{1}{2}$ gilt: $\mu_S(n)=0$ und $\sigma_S^2(n)=n\delta^2$.


\section{Markovketten und \\ bedingte Unabhangigkeit}

\CheatsheetHeading{Markov-Ungleichung:}{}

Abschatzung fur große Ausschlage einer ZV: \\
$\mathbb{P}(\{|X|\ge a\}) \le \frac{\mathbb{E}[|X|]}{a},\quad a>0$

\CheatsheetHeading{Tschebyschow-Ungleichung:}{}

$\mathbb{P}(\{|X-\mathbb{E}[X]|\ge a\}) \le \frac{\operatorname{Var}[X]}{a^2},\quad a>0$ \\
$\Rightarrow \mathbb{P}(\{|X-\mathbb{E}[X]|<a\}) \ge 1-\frac{\operatorname{Var}[X]}{a^2}$

\CheatsheetHeading{Bedingte Unabhangigkeit:}{}

Ereignisse $A,C$ sind bedingt unabhangig bzgl. $B$, wenn \\
$\mathbb{P}(A\cap C\mid B)=\mathbb{P}(A\mid B)\mathbb{P}(C\mid B)$ \\
$\Leftrightarrow \mathbb{P}(A\mid B\cap C)=\mathbb{P}(A\mid B)$ \\
$\Leftrightarrow \mathbb{P}(C\mid B\cap A)=\mathbb{P}(C\mid B)$ \\
\textbf{Achtung:} i.A. gilt nicht $\mathbb{P}(A\cap C)=\mathbb{P}(A)\mathbb{P}(C)$.

{\color{orange} \textbf{Fur Zufallsvariablen/pdf/pmf:}}

$p_{X,Z\mid Y}(x,z\mid y)=p_{X\mid Y}(x\mid y)\,p_{Z\mid Y}(z\mid y)$ \\
$p_{Z\mid Y,X}(z\mid y,x)=p_{Z\mid Y}(z\mid y)$ \\
$f_{Z\mid Y,X}(z\mid y,x)=f_{Z\mid Y}(z\mid y)$ \\
$f_{X,Z\mid Y}(x,z\mid y)
    = f_{X\mid Y}(x\mid y)\cdot f_{Z\mid Y}(z\mid y)$ \\ [3pt]
$\Rightarrow\quad\frac{f_{X_1,X_2,X_3}(x_1,x_2,x_3)}{f_{X_1,X_2}(x_1,x_2)}
    = \frac{f_{X_2,X_3}(x_2,x_3)}{f_{X_2}(x_2)}$ \\ [-3pt] 

$\Rightarrow\quad
\frac{f_{X_1,X_2,X_3}(x_1,x_2,x_3)}{f_{X_2,X_3}(x_2,x_3)}
    = \frac{f_{X_1,X_2}(x_1,x_2)}{f_{X_2}(x_2)}$ \\ [-5pt]


\CheatsheetHeading{Markov-Eigenschaft:}{}

$\mathbb{P}(X_{n+1}=x_{n+1}\mid X_n,\dots,X_1)=\mathbb{P}(X_{n+1}=x_{n+1}\mid X_n)$ \\
Beispiel (1. Ordnung): $f_{X_3\mid X_2,X_1}(x_3\mid x_2,x_1)=f_{X_3\mid X_2}(x_3\mid x_2)$

\CheatsheetHeading{Homogene Markovkette:}{}

$p_{ij}=\mathbb{P}(X_{n+1}=j\mid X_n=i)$ ist unabhangig vom Index $n$.

\CheatsheetHeading{Zustandsubergang / Matrix:}{}

Ubergangsmatrix: $\Pi=[p_{ij}]$,\ $p_{ij}\in[0,1]$,\ $\sum_j p_{ij}=1$ \\
$m$-Schritt-Ubergange: $p_{ij}^{(m)}=\mathbb{P}(X_{n+m}=j\mid X_n=i)=[\Pi^m]_{ij}$

\CheatsheetHeading{Zustandsvektor:}{}

$\underline{p}_{n+1}=\underline{p}_n\Pi,\quad \underline{p}_{n+m}=\underline{p}_n\Pi^m$

\CheatsheetHeading{Stationare Verteilung:}{}

$\underline{\pi}=\underline{\pi}\Pi,\quad \sum_i \pi_i=1$ \\
Bei endlichem, irreduziblem und aperiodischem Zustandraum ist $\underline{\pi}$ eindeutig und
$\underline{p}_n\to\underline{\pi}$.


\section{Reelle Zufallsprozesse}

\CheatsheetHeading{Grundgroßen:}{}

$\mu_X(t)=\mathbb{E}[X_t]$ \\
$\sigma_X^2(t)=\operatorname{Var}[X_t]$ \\
$r_X(t,s)=\mathbb{E}[X_tX_s]$ \\
$c_X(t,s)=\operatorname{Cov}[X_t,X_s]=r_X(t,s)-\mu_X(t)\mu_X(s)$

\CheatsheetHeading{Gemeinsame Prozesse $X_t,Y_t$:}{}

$r_{XY}(t,s)=\mathbb{E}[X_tY_s]$ \\
$c_{XY}(t,s)=r_{XY}(t,s)-\mu_X(t)\mu_Y(s)$

\CheatsheetHeading{WSS-Prozess:}{}

$\mu_X(t)=\mu_X$ konstant,\quad $r_X(t,s)=r_X(t-s)$,\quad $c_X(t,s)=c_X(t-s)$

\CheatsheetHeading{Orthogonalitat / Unkorreliertheit:}{}

Orthogonal: $r_{XY}(t,s)=0$ \\
Unkorreliert: $c_{XY}(t,s)=0$

\CheatsheetHeading{Wiener-Prozess (kurz):}{}

$W_t\sim\mathcal{N}(0,\sigma^2 t)$,\quad $\mathbb{E}[W_t]=0$,\quad $\operatorname{Var}[W_t]=\sigma^2 t$, \\
$r_W(s,t)=\sigma^2\min(s,t)$ und unabhangige Inkremente.


\section{Poisson Prozess}

\CheatsheetHeading{Definition:}{}

$N_t$: Anzahl Ereignisse in $[0,t]$ mit Rate $\lambda>0$.

$N_t\sim\operatorname{Pois}(\lambda t)$ \\
$\mathbb{P}(N_t=n)=\frac{(\lambda t)^n}{n!}\mathrm{e}^{-\lambda t},\ n\in\mathbb{N}_0$

\CheatsheetItemListStart
  \CheatsheetItem{Eigenschaften:}{}{}
\CheatsheetItemListEnd
$\sbullet \mathbb{E}[N_t]=\lambda t$ \\
$\sbullet \operatorname{Var}[N_t]=\lambda t$ \\
$\sbullet r_N(s,t)=\lambda\min(s,t)+\lambda^2st$ \\
$\sbullet c_N(s,t)=\lambda\min(s,t)$ \\
$\sbullet Inkremente sind unabhangig und stationar$ \\
$\sbullet\ X_i\sim\operatorname{Exp}(\lambda)\ \text{(Wartezeiten, i.i.d.)}$


\section{Leistungsdichtespektrum und LTI-Systeme}

\CheatsheetHeading{LTI-System:}{}

$y(t)=(h*x)(t)=\int_{-\infty}^{\infty}h(\tau)x(t-\tau)\,d\tau$ \\
$Y(f)=H(f)X(f)$

\CheatsheetHeading{Bei WSS-Eingang:}{}

$\mu_Y = \mu_X \int h(t)\,dt = \mu_XH(0)$ \\
$r_{YY}(\tau) = (h*\tilde h*r_{XX})(\tau),\ \tilde h(t)=h(-t)$

\CheatsheetHeading{Leistungsdichte (LDS):}{}

$S_{YY}(f)=|H(f)|^2S_{XX}(f)$ \\
$S_{XY}(f)=H(f)S_{XX}(f),\quad S_{YX}(f)=H^*(f)S_{XX}(f)$

\CheatsheetHeading{Allgemein fur WSS-Prozesse:}{}

$\sbullet S_X(f)\ge 0$ \\
$\sbullet S_X(-f)=S_X(f)$ fur reelle Prozesse \\
$\sbullet \int_{-\infty}^{\infty} S_X(f)\,df = \operatorname{Var}[X] + \mu_X^2$







%-------------------------------------------
%%%%%%  CHEATSHEET ENDS HERE  %%%%%%%%%%%%
%-------------------------------------------

\end{multicols*}
\end{document}
